\documentclass[11pt]{article}
\usepackage{icmlko}
\begin{document}

\title{결측으로 버는 축을 세어 봤다 --- 서른셋 중 하나, 그리고 그 하나도 진짜였다}
\author{월드모델 실험실}
\date{노트 355 · 2026년 8월 1일}
\maketitle

\begin{abstract}
두 번 걸렸으면 세어 봐야 한다. 노트 266이 애니 태그 성분에서, 노트 354가
시장팝업 \texttt{entry\_friction} 에서 ``축이 버는 것은 값이 아니라
결측''을 잡았다. 판의 모든 (도메인, 축) 쌍에 같은 검사를 돌렸다 --- 값은
\emph{관측된 행 안에서만}, 결측은 ``적혔나''로, 둘 다 연도 통제.
\textbf{잴 수 있는 서른세 쌍 중 깃발은 하나}(아이돌
\texttt{wiki\_momentum})이고, 그 하나도 파 보니 진짜였다: 사후 표시자라면
최근 데뷔가 더 벌어야 하는데 \emph{옛} 데뷔가 $+0.4204$, 최근이
$+0.2411$ 로 반대다. 문서 있음/없음 격차는 다섯 연도 구간 전부에 있고,
위키 계열을 아이돌에서 빼면 \textbf{아이돌 $-0.0668$}(1/12)이다.
\textbf{판은 결측을 타고 있지 않다} --- 두 번은 두 번이었지 무늬가 아니었다.
덤으로 둘. ① \textbf{표시자는 축이 아니라 \emph{계열}의 성질이다} --- 위키
셋과 검색 셋이 표시자 하나를 나눠 쓴다. ② 노트 354의 그 축은 \emph{이제
안 걸린다}. MKT2 행이 결측인데 라벨은 평범해서 결측 신호가
$+0.2362 \to +0.0224$ 로 흩어졌다 --- 고침이 축을 막았을 뿐 아니라 그
신호를 없앴다.
\end{abstract}

\section{왜 세나}

노트 266: 애니 \texttt{air\_quarter} 태그 성분 넷이 검사를 통과했는데
기록된 행 안에서 다시 재니 전부 $|\rho|<0.07$ 이었다 --- 신호가 통째로
결측 무늬였다. 노트 354: 시장팝업 \texttt{entry\_friction} 의 값은 유보
$+0.0664$ 인데 ``적혔나''는 $+0.2362$ 였고, 그것이 그 도메인 $\rho$ 의
절반을 대고 있었다.

두 번이면 묻게 된다 --- \emph{판 전체가 이러고 있나?} 축이 서른아홉인데
아무도 센 적이 없다. 셌다.

\textbf{검사.} (도메인, 축)마다 연도 통제 스피어만 둘. ① \textbf{값} ---
표시자가 켜진 행 안에서만 값과 라벨. ② \textbf{결측} --- 전체 행에서
``관측됐나''와 라벨. 모형 적합이 없다. 깃발은 미리 정했다: 표시자 평균이
$0.05$\textasciitilde$0.95$, 학습 관측 30행 이상, $|$결측$| \ge 0.10$,
$|$결측$| \ge 2|$값$|$, 그리고 \textbf{학습과 유보 둘 다}에서.

\section{서른셋 중 하나}

\begin{center}\includegraphics[width=\linewidth]{figs/mask.pdf}\end{center}

잴 수 있는 쌍이 서른셋이고(나머지는 표시자가 상수라 잴 게 없다) 깃발은
\textbf{하나} --- 아이돌 \texttt{wiki\_momentum} 이다(값 유보 $-0.1592$ ·
결측 유보 $+0.3495$).

\textbf{먼저 구조를 하나 배웠다: 표시자는 축이 아니라 계열의 성질이다.}
표에서 아이돌 \texttt{wiki\_level} · \texttt{wiki\_momentum} ·
\texttt{wiki\_volatility} 의 결측 값이 \emph{똑같이} $+0.3495$ 다. 셋이
같은 관측 표시자를 나눠 쓰기 때문이다. 검색 셋도 마찬가지다. 그러니
``어느 \emph{축}이 결측으로 버나''는 잘못 물은 것이고, 옳은 물음은
``이 \emph{계열}의 값이 제 표시자를 데리고 다닐 만한가''다.
\texttt{wiki\_momentum} 이 혼자 걸린 것은 그 축이 죽어 있어서지
(값 유보 $-0.1592$) 계열이 나빠서가 아니다 --- 같은 표시자를 쓰는
\texttt{wiki\_level} 은 값 유보 $+0.5166$ 이다.

\section{그 하나를 판다}

노트 306이 ``위키 표시자가 사후일 수 있다''고 적어 뒀다 --- 2026년에
문서가 있다는 것은 그 사이에 떴다는 뜻일 수 있다. 사전 등록에서도 위키를
제일 의심했다. 셋으로 갈랐다.

\textbf{하나 --- 사후라면 최근이 더 벌어야 한다.} 2015년 데뷔 팀은 문서가
생길 시간이 십 년 있었지만 2026년 데뷔 팀에 몇 달 만에 문서가 있다면
그것은 성공의 결과일 확률이 높다. 그러면 표시자$\leftrightarrow$라벨이
최근에서 더 세야 한다.

\begin{center}
\begin{tabular}{lrrr}
\hline
데뷔 & $n$ & 덮음 & 표시자 $\leftrightarrow$ 라벨 \\
\hline
2023년까지 & 99 & 0.44 & $\mathbf{+0.4204}$ \\
2024년부터 & 74 & 0.49 & $+0.2411$ \\
\hline
\end{tabular}
\end{center}

\textbf{반대다.} 옛 데뷔가 더 세다.

\textbf{둘 --- 격차가 시기를 안 탄다.} 문서 있음/없음의 초동 중앙 차가
$+1.33 \cdot +0.88 \cdot +0.40 \cdot +0.84 \cdot +0.52$ 로 다섯 구간에
전부 있다. 덮음은 $0.47 \cdot 0.36 \cdot 0.64 \cdot 0.55 \cdot 0.20$ 으로
최근일수록 높지도 않다.

\textbf{셋 --- 빼 보면 손해다.} 위키 계열을 아이돌에서 빼고 열두 번 짝
뽑기: 판 $-0.0007$(3/12) · \textbf{아이돌 $-0.0668$(1/12)}.

셋이 같은 말을 한다. \textbf{위키 문서가 있다는 것은 실제로 큰 팀이라는
뜻이다} --- 세계의 성질이지 우리가 만든 자국이 아니다. 사전 등록에
``깃발은 조사 대상이지 판정이 아니다''라고 적어 둔 그대로다. 안 막는다.

\section{덤 둘}

\textbf{하나 --- 노트 354의 그 축은 이제 안 걸린다.} 시장팝업
\texttt{entry\_friction} 의 결측 신호가 유보 $+0.2362 \to \mathbf{+0.0224}$
로 흩어졌다. MKT2 44행이 표시자 0 인데 라벨은 평범해서
(중앙 3.23 대 3.14), ``표시자 0''이 더 이상 ``안 적힐 만큼 작은 행사''를
뜻하지 않는다. \textbf{노트 354의 고침은 축을 막았을 뿐 아니라 그 신호가
서 있던 땅을 없앴다.} 노트 354의 $+0.2362$ 는 MKT 행만으로 잰 것이고 그
글에 그렇게 적혀 있다 --- 두 수는 다른 집합의 수다.

\textbf{둘 --- 덮음이 시기를 타면 가짜 결측 신호가 생긴다.} 도서 검색
계열은 결측이 학습 $-0.0002$ 인데 유보 $\mathbf{-0.2554}$ 다. 덮음이
$0.69 \to 0.99$ 로 뛰어서다 --- 유보에서 표시자 0 인 것은 한 줌뿐이고
그 한 줌이 우연히 낮다. 미리 정한 ``학습과 유보 \textbf{둘 다}''가
이것을 걸러 냈다. \textbf{한쪽만 보고 깃발을 세웠으면 도서 검색을 막을
뻔했다.}

\section{그래서}

\begin{itemize}
\item \textbf{판은 결측을 타고 있지 않다.} 서른세 쌍 중 걸린 하나도
파 보니 진짜였다. 노트 266과 354는 두 번이었지 무늬가 아니었다.
\item \textbf{검사를 계열 단위로 다시 정의한다.} 표시자는 계열이 공유한다.
물어야 할 것은 ``이 계열의 값들이 제 표시자를 데리고 다닐 만한가''다.
\item \textbf{덮음이 시기를 타는 계열은 따로 표시한다} --- 도서 검색
$0.69 \to 0.99$. 결측 신호가 유보에만 나타나면 축이 아니라 수집 이력이다.
\item \textbf{안 막은 것도 적는다.} 아이돌 위키는 깃발이 섰고 살렸다.
근거 셋을 위에 적어 뒀다 --- 다음에 누가 다시 물으면 이 문단이 답이다.
\end{itemize}

\section{남은 것}

\begin{center}
\begin{tabular}{ll}
\hline
갈래 & 지금 아는 것 \\
\hline
\texttt{wiki\_momentum} 셋 & 아이돌 · 모바일 · 세계애니에서 값이 죽어 \\
 & 있다(유보 $-0.159 \cdot -0.236 \cdot +0.039$). 계열은 살리고 \\
 & 이 축만 뺄 수 있나 --- 다만 판 가중으로는 SD 아래다 \\
MKT2 나머지 195건 & 라벨 · 기간 없음. 다시 뽑으면 유보가 더 는다 \\
MKT2 입장/예약 0/44 & 원문에서 다시 뽑으면 축이 되살아난다 \\
추가 2026 아홉 행 & 노트 353. 아이돌 라벨 수준이 한터와 반대 \\
\hline
\end{tabular}
\end{center}

첫 줄은 노트 351의 환율로 먼저 계산했다 --- 아이돌 $51/2{,}717 = 1.9\%$,
모바일 $441/2{,}717 = 16.2\%$. 모바일 \texttt{wiki\_momentum} 은 덮음이
$0.05$ 라 계열이 거의 안 붙어 있어 움직일 것이 없고, 아이돌은 가중이
얇다. \textbf{계산이 ``돌리지 마라''고 한다.} 노트 351이 사 준 것이
이것이다.

\end{document}
